\documentclass[11pt,a4paper]{article}
\usepackage[review]{acl}
\usepackage{times}
\usepackage{latexsym}
\usepackage{graphicx}
\usepackage{booktabs}
\usepackage{amsmath}

\title{A Dual-Reference Metric for Dialect Preservation in {ASR}: \\ Evidence from {German} Dialect Continua}

\author{Anonymous ACL submission}

\begin{document}
\maketitle

\begin{abstract}
Word Error Rate (WER) evaluates ASR against a single reference---a premise that fails for dialect continua like Swiss German, where no standardized orthography exists. Prior work has diagnosed this mismatch (FlexWER, PolyWER) and quantified its magnitude (Akeret finds 44--46\% of measured WER on Swiss German is stylistic variation, not error). Building on these observations, we operationalize the \textbf{Dialect Preference Score (DPS)}, the difference between normalized character similarity to a parallel dialect reference and to a standard language reference. DPS adds a dimension WER is blind to: where an ASR output falls on the dialect--standard continuum. Across eight models from three families on Swiss and German dialect data, we find three stable behavioral patterns: CTC models lean dialectal, Whisper models drift toward standard German with increasing scale (from $+0.04$ to $-0.12$), and LLM-ASR models occupy a high-similarity equilibrium near zero. Causal experiments confirm the decoder LM prior drives standardization, with effect sizes growing monotonically with model capacity. Component ablation shows neither dialect-similarity nor standard-similarity alone captures these patterns.
\end{abstract}

\section{Introduction}

WER compares an ASR hypothesis against \textit{one} reference---a premise rarely questioned: each utterance has exactly one correct written form. For dialect continua, this premise is a category error. Swiss German lacks standardized orthography; professional subtitlers at Swiss public television produce Standard German text from Swiss German speech as a matter of policy, acknowledging the task is speech-to-translation, not transcription. WER penalizes genuine misrecognition and valid dialect--standard paraphrase identically. Language models more broadly exhibit a ``default to standard'' bias with non-standard varieties~\shortcite{fleisig2024linguistic}, and dialectal evaluation disparities are documented across text NLP~\shortcite{deas2023evaluation}.

This is not a new observation. Nigmatulina et al.~\shortcite{Nigmatulina2020ASRFN} proposed FlexWER, using normalization dictionaries to relax exact matching. Kadaoui et al.~\shortcite{Kadaoui2024PolyWERAH} introduced PolyWER, accepting multiple reference forms but aggregating to a single score. Akeret~\shortcite{Akeret2026SubtitleAlignedFO} quantified the magnitude: 44--46\% of measured WER on Swiss German is valid stylistic variation. Blaschke et al.~\shortcite{Blaschke2025StandardtoDialectTT} showed that computing separate WER against dialect and standard references reveals systematic normalization in Whisper.

These contributions share a diagnosis: single-reference WER is structurally inadequate for dialect speech. What has been missing is a simple way to \textit{position} output on the dialect--standard continuum---not how close it is to one reference, but \textit{which side it leans toward}. This paper provides that instrument.

We define DPS as $\text{DialectSim} - \text{NormSim}$: character similarity to a dialect reference minus similarity to a standard reference. Positive values indicate dialect-leaning output; negative values indicate standardization drift. It requires no manual rules or training, only parallel transcriptions---a resource provided by SwissDial~\shortcite{dogan2021swissdial} and Betthupferl.~\shortcite{Blaschke2025AMD}

Morphologically complex, non-standardized varieties are known to challenge language modeling~\shortcite{cotterell2018all}. We study German dialect continua precisely \textit{because} their orthography is maximally unstandardized. If a dual-reference metric cannot produce interpretable signal under these conditions, it is unlikely to matter anywhere. Conversely, if it does, the patterns it reveals should generalize: the decoder LM prior is a property of the model, not of the language. Across eight models from three families, we document:

\begin{enumerate}
\item \textbf{Architecture signatures.} CTC: positive DPS (dialect-leaning). Whisper: monotonically decreasing with scale ($+0.04$ to $-0.12$). LLM-ASR: near-zero equilibrium despite highest absolute similarity.
\item \textbf{Causal mechanism.} Prompt calibration, temperature sweeps, and long-form tests converge on the decoder LM prior as the driver of standardization, with effect sizes growing with model capacity.
\item \textbf{Component necessity.} DialectSim alone would rank Qwen3-ASR as the best dialect preserver---a conclusion DPS reveals as spurious.
\end{enumerate}

\section{DPS: A Dual-Reference Diagnostic}

Let $y$ be an ASR hypothesis, $r_d$ a dialect reference, and $r_s$ a standard reference:
\begin{align}
\text{DialectSim}(y, r_d) &= 1 - \frac{\text{NED}(y, r_d)}{\max(|y|, |r_d|)} \label{eq:ds} \\
\text{NormSim}(y, r_s) &= 1 - \frac{\text{NED}(y, r_s)}{\max(|y|, |r_s|)} \label{eq:ns} \\
\text{DPS} &= \text{DialectSim} - \text{NormSim} \label{eq:dps}
\end{align}
where NED is normalized character edit distance. We use character-level similarity because dialect orthography lacks consistent word boundaries, making token-level metrics fragile to segmentation variance. Both components range in $[0,1]$; DPS in $[-1, 1]$. When $\max(\text{DialectSim}, \text{NormSim}) < 0.1$, the output is classified as model collapse and excluded.

\textbf{Why the difference, not two scores.} Reporting separate WER against each reference tells you distance to each anchor; DPS tells you the \textit{relationship} between them. Output that is 0.70 similar to both references (DPS $= 0.00$) is qualitatively different from one that is 0.80 to dialect and 0.60 to standard (DPS $= +0.20$), despite equal average similarity. The sign captures directional preference; the magnitude captures its strength.

\section{Experimental Setup}

\textbf{Datasets.} SwissDial~\shortcite{dogan2021swissdial} provides $\sim$24h of read speech across 8 Swiss German dialects with parallel dialect and Standard German transcriptions. Betthupferl~\shortcite{Blaschke2025AMD} provides $\sim$4h across 7 Southern German dialect regions (Bavarian, Franconian, Alemannic) from 70 speakers, also with paired transcriptions. We use SwissDial as our primary benchmark and Betthupferl for cross-dataset validation.

\textbf{Models.} We evaluate eight models zero-shot across three families: (1) \textbf{CTC}: wav2vec2-XLS-R 300M and 1B,~\shortcite{Babu2021XLSRSC} lacking an autoregressive decoder; (2) \textbf{Encoder-decoder}: Whisper tiny, small, medium, large-v3,~\shortcite{Radford2022RobustSR} with increasing decoder LM prior; (3) \textbf{LLM-ASR}: Qwen3-ASR 0.6B and 1.7B, whose decoder is a full language model backbone.

\textbf{Causal interventions.} To test whether the decoder LM prior drives DPS differences, we apply three inference-time interventions (no weight updates): (i) \textit{Prompt calibration}: following Yang et al.~\shortcite{yang2024prompts}, who showed Whisper's decoder is sensitive to prompt linguistic variety but not its semantic content, we transcribe each utterance with empty, dialect German, and standard German prompts sharing identical semantics. (ii) \textit{Temperature sweep}: $T \in \{0.0, 0.2, 0.5, 0.8, 1.0\}$; higher $T$ dilutes the LM prior. (iii) \textit{Long-form toggle}: concatenate 500 same-dialect sentences into $\sim$40min audio spanning $\sim$80 Whisper chunks, then compare $\texttt{condition\_on\_previous\_text}=\text{True}$ vs.\ $\text{False}$.

\section{Results}

\subsection{Architecture Signatures}

Table~\ref{tab:arch} reports DPS and components. Three behavioral patterns emerge, stable across all eight dialects.

\begin{table}[t]
\centering\small
\setlength{\tabcolsep}{4pt}
\begin{tabular}{lccc}
\toprule
\textbf{Model} & \textbf{DialectSim} & \textbf{NormSim} & \textbf{DPS} \\
\midrule
\multicolumn{4}{c}{\textit{CTC (acoustic-only)}} \\
wav2vec2 300M & .670 & .567 & \textbf{+.103} \\
wav2vec2 1B   & .673 & .594 & \textbf{+.079} \\
\midrule
\multicolumn{4}{c}{\textit{Whisper (encoder-decoder)}} \\
tiny   & .647 & .608 & \textbf{+.039} \\
small  & .699 & .753 & \textbf{--.055} \\
medium & .712 & .816 & \textbf{--.104} \\
large-v3 & .718 & .834 & \textbf{--.116} \\
\midrule
\multicolumn{4}{c}{\textit{LLM-ASR (Qwen3)}} \\
0.6B & .728 & .727 & \textbf{+.001} \\
1.7B & .747 & .764 & \textbf{--.016} \\
\bottomrule
\end{tabular}
\caption{DPS across three model families. All values are bootstrap means ($n=2{,}000$ resamples). 95\% CIs are tight ($\pm 0.001$--$0.002$) and non-overlapping across families. CTC: dialect-leaning. Whisper: monotonically decreasing with scale. LLM-ASR: near-zero equilibrium.}
\label{tab:arch}
\end{table}

\textbf{CTC: acoustic faithfulness.} Without a decoder LM, wav2vec2 models produce positive DPS ($+0.08$--$+0.10$). Their directional behavior is distinct: they transcribe what they hear.

\textbf{Whisper: capacity-driven standardization.} DPS decreases monotonically from $+0.039$ (tiny) to $-0.116$ (large-v3). Large-v3's DialectSim (0.718) exceeds wav2vec2 1B's (0.673), yet its NormSim gain (0.834 vs.\ 0.594) far outpaces its DialectSim gain, pulling DPS negative: the model overrides acoustic evidence with a standard-language prior that strengthens with scale.

\textbf{LLM-ASR: high-similarity equilibrium.} Qwen3-ASR achieves the highest scores on \textit{both} axes among all tested models, yielding DPS near zero ($+0.001$ and $-0.016$). This is not neutral behavior. It reflects a third pattern distinct from both CTC and Whisper. The model's full language model backbone enables simultaneous high fidelity to both dialectal and standard forms, producing a hybrid register that is syntactically dialect-like but orthographically standard-leaning. We hypothesize tokenizer granularity contributes: Whisper's $\sim$52K-token BPE vocabulary sparsely represents dialect variants, while Qwen3-ASR inherits $\sim$152K tokens from its Qwen3 backbone. Training data composition differences, not publicly documented in comparable detail, further complicate attribution. DPS makes this hybrid pattern visible; isolating causes requires controlled experiments beyond this study's scope.

\subsection{Causal Attribution}

\textbf{Prompt calibration.} Table~\ref{tab:prompt} shows DPS under three prompt conditions. A 30-word dialect German prompt shifts DPS toward dialect across all scales, with the effect growing from $+0.015$ (tiny) to $+0.139$ (large-v3). Medium Whisper reverses sign entirely (from $-0.098$ empty to $+0.053$ dialect). The standard German prompt barely differs from empty ($|\Delta| < 0.01$), confirming Whisper's default is standard behavior.

\begin{table}[t]
\centering\small
\begin{tabular}{lccc|c}
\toprule
\textbf{Model} & \textbf{Empty} & \textbf{Dialect} & \textbf{Standard} & $\bm{\Delta}$ \\
\midrule
tiny   & +.033 & +.043 & +.028 & +.015 \\
small  & --.069 & --.009 & --.067 & +.058 \\
medium & --.098 & +.053 & --.094 & +.146 \\
large-v3 & --.129 & +.018 & --.121 & +.139 \\
\bottomrule
\end{tabular}
\caption{Prompt calibration. $\Delta = \text{DPS}(\text{dialect}) - \text{DPS}(\text{standard})$, all bootstrap 95\% CIs within $\pm 0.003$. Effect grows monotonically with model capacity; all pairwise CIs are non-overlapping.}
\label{tab:prompt}
\end{table}

\textbf{Temperature sweep.} The DPS difference between $T=1.0$ and $T=0.0$ is $+0.004$ (small, 95\% CI $[+0.002,+0.006]$), $+0.017$ (medium, 95\% CI $[+0.014,+0.019]$), and $+0.015$ (large-v3, 95\% CI $[+0.012,+0.018]$). The effect is smaller than prompt calibration (temperature is a blunt instrument), but directionally consistent and statistically distinguishable from zero for medium and large-v3.

\textbf{Long-form cross-chunk feedback.} In single utterances (3--7s, one chunk), $\texttt{condition\_on\_previous\_text}$ has no measurable effect ($|\Delta| < 0.002$). In the concatenated setting ($\sim$80 chunks), large-v3 shows $\text{DPS}(\text{False}) - \text{DPS}(\text{True}) = +0.024$ (95\% CI~$[+0.016,+0.029]$) with 8/8 dialects agreeing on direction. Smaller models show no systematic effect, consistent with their weaker LM priors.

\subsection{Component Ablation}

Table~\ref{tab:ablation} shows what each component of DPS would conclude independently.

\begin{table}[t]
\centering\small
\begin{tabular}{lccc}
\toprule
\textbf{Claim} & \textbf{DialectSim} & \textbf{NormSim} & \textbf{DPS} \\
\midrule
``CTC preserves dialect'' & $\times$ (low abs.) & $\times$ & \checkmark \\
``Whisper: bigger = more standard'' & $\times$ (reverse) & \checkmark & \checkmark \\
``Qwen3-ASR is best at dialect'' & \checkmark & \checkmark & $\times$ \\
\midrule
Produces architecture signatures? & No & No & Yes \\
Responds directionally? & Partially & Partially & Yes \\
Reveals \textit{what} changed? & No & No & Yes \\
\bottomrule
\end{tabular}
\caption{Component ablation. $\checkmark$ = supports claim; $\times$ = contradicts or misleads. Only DPS produces the tripartite architecture pattern and clean directional responses.}
\label{tab:ablation}
\end{table}

DialectSim alone ranks Qwen3-ASR 0.6B (0.728) and large-v3 (0.718) as the top two ``dialect preservers.'' DPS reveals this ranking as an artifact of overall fluency: Qwen3-ASR's NormSim is equally high (0.727), and large-v3 achieves high DialectSim through strong general language modeling rather than dialect-specific fidelity. A researcher using DialectSim alone would conclude the largest models are most dialect-preserving---the opposite of what DPS shows. NormSim alone tracks conventional WER, capturing the Whisper scale trend but unable to distinguish improved accuracy from dialect erasure. The signed difference operationalizes a shared insight: dialect evaluation requires going beyond proximity to a single reference, toward capturing the \textit{tension} between competing transcription norms.

\section{Discussion}

We have shown that a simple dual-reference metric reveals systematic behavioral differences across ASR architectures invisible to single-reference evaluation.

The three architecture signatures (CTC positive, Whisper monotonically negative, LLM-ASR near-zero) are stable across all eight Swiss German dialects and replicate on Betthupferl, where Whisper DPS decreases from $+0.002$ (tiny) to $-0.088$ (large-v3), reproducing the same monotonic trend on independent data. Reporting two separate WER scores (as in Blaschke et al.) would show large-v3 with the lowest WER against \textit{both} references---suggesting ``best overall performance.'' DPS reveals this comes at the cost of the strongest standardization drift, a dimension invisible to separate reporting.

The causal evidence reinforces this picture. Three independent manipulations of the decoder's text prior (prompt calibration, temperature, long-form) all shift DPS in the predicted direction, with effect sizes growing with model capacity. Prompt calibration is particularly informative: a 30-word dialect sentence as initial prompt reverses the DPS sign for medium and large-v3, directly demonstrating the decoder LM prior drives standardization.

The component ablation confirms DPS captures something neither component captures alone: DialectSim misranks models, NormSim is blind to dialect, and only the signed difference produces interpretable architecture signatures.

\textbf{What DPS does and does not claim.} DPS does not replace WER. It adds a dimension WER is structurally blind to. For varieties without a fixed written standard, single-reference metrics answer the wrong question. DPS answers a different one: does this model lean dialect or standard, and by how much? The answer can inform model selection, fairness auditing of deployed speech systems, and inference-time dialect preservation strategies.

\textbf{Limitations.} (i) Behavioral signatures replicate across SwissDial and Betthupferl (15 regional varieties); causal experiments are restricted to SwissDial due to Betthupferl audio access constraints. We release inference and analysis code to enable replication on any dataset with parallel transcriptions. (ii) Character-level similarity captures orthographic and phonetic variation but may miss morphosyntactic dialect features (e.g., word order, tense choice). Token-level and embedding-based variants may capture morphological features that character-level similarity misses. (iii) We validated DPS through architecture consistency, causal responsiveness, and cross-dataset replication, but not through direct human correlation; recruiting dialect speakers for perceptual validation remains an important next step. (iv) The causal experimental paradigm currently covers the Whisper family. Prompt calibration and temperature sweeps can extend to LLM-ASR models; the long-form toggle has no direct equivalent, as Qwen3-ASR processes audio as a single input without Whisper's chunked architecture. (v) DPS requires parallel dialect--standard transcriptions. While such resources remain scarce, the growing availability of dual-reference datasets for dialectal speech suggests this constraint will ease over time.

These results point toward a recalibration of how the field evaluates ASR for non-standard speech varieties. The question should not be ``how accurate is this transcription?'' alone. It should also be ``what kind of language is this system producing?'' DPS offers one answer to the second question. It is not the only possible answer, but it is a starting point: simple to compute, grounded in existing dataset infrastructure, and capable of revealing architectural patterns that single-reference evaluation leaves invisible.

\section*{Acknowledgments}
We thank the creators of SwissDial and Betthupferl for making their data publicly available.

\bibliographystyle{acl_natbib}
\bibliography{references}

\end{document}
